%
% Main document
% ===========================================================================
% This is part of the document "Project documentation template".
% Authors: brd3, kaa1
%

%---------------------------------------------------------------------------
\documentclass[
	a4paper,					% paper format
	10pt,							% fontsize
	twoside,					% double-sided
	notitlepage,			% use no standard title page
	parskip=half,			% set paragraph skip to half of a line
]{scrreprt}					% KOMA-script report
%---------------------------------------------------------------------------

\raggedbottom
\KOMAoptions{cleardoublepage=plain}			% Add header and footer on blank pages


% Load Standard Packages:
%---------------------------------------------------------------------------
\usepackage[standard-baselineskips]{cmbright}
\usepackage{svg}
\usepackage{underscore}
\usepackage{transparent}
\usepackage{import}
\usepackage[english]{babel}										% german hyphenation
\usepackage[utf8]{inputenc}  							% Unix/Linux - load extended character set (ISO 8859-1)
%\usepackage[ansinew]{inputenc}  							% Windows - load extended character set (ISO 8859-1)
\usepackage[T1]{fontenc}											% hyphenation of words with ä,ö and ü
\usepackage{textcomp}													% additional symbols
\usepackage{ae}																% better resolution of Type1-Fonts
\usepackage{fancyhdr}													% simple manipulation of header and footer
\usepackage{etoolbox}													% color manipulation of header and footer
\usepackage{graphicx}                      		% integration of images
\usepackage{float}														% floating objects
\usepackage{caption}													% for captions of figures and tables
\usepackage{booktabs}													% package for nicer tables
\usepackage{tocvsec2}													% provides means of controlling the sectional numbering
\usepackage{listings}
%---------------------------------------------------------------------------

% Load Math Packages
%---------------------------------------------------------------------------
\usepackage{amsmath}                    	   	% various features to facilitate writing math formulas
\usepackage{amsthm}                       	 	% enhanced version of latex's newtheorem
\usepackage{amsfonts}                      		% set of miscellaneous TeX fonts that augment the standard CM
\usepackage{amssymb}													% mathematical special characters
\usepackage{exscale}													% mathematical size corresponds to textsize
%---------------------------------------------------------------------------

% Package to facilitate placement of boxes at absolute positions
%---------------------------------------------------------------------------
\usepackage[absolute]{textpos}
\setlength{\TPHorizModule}{1mm}
\setlength{\TPVertModule}{1mm}
%---------------------------------------------------------------------------

% Definition of Colors
%---------------------------------------------------------------------------
\RequirePackage{color}                          % Color (not xcolor!)
\definecolor{linkblue}{rgb}{0,0,0.8}            % Standard
\definecolor{darkblue}{rgb}{0,0.08,0.45}        % Dark blue
\definecolor{bfhgrey}{rgb}{0.41,0.49,0.57}      % BFH grey
%\definecolor{linkcolor}{rgb}{0,0,0.8}     			% Blue for the web- and cd-version!
\definecolor{linkcolor}{rgb}{0,0,0}        			% Black for the print-version!
%---------------------------------------------------------------------------

% Hyperref Package (Create links in a pdf)
%---------------------------------------------------------------------------
\usepackage[
	pdftex,english,bookmarks,plainpages=false,pdfpagelabels,
	backref = {false},										% No index backreference
	colorlinks = {true},                  % Color links in a PDF
	hypertexnames = {true},               % no failures "same page(i)"
	bookmarksopen = {true},               % opens the bar on the left side
	bookmarksopenlevel = {0},             % depth of opened bookmarks
  pdftitle = {{Presentator3000, an online presentation tool}},	   	% PDF-property
  pdfauthor = {{bollm6, schmidrobin}},        					  % PDF-property
	pdfsubject = {Presentator3000},        % PDF-property
	linkcolor = {linkcolor},              % Color of Links
	citecolor = {linkcolor},              % Color of Cite-Links
	urlcolor = {linkcolor},               % Color of URLs
]{hyperref}
%---------------------------------------------------------------------------

% Set up page dimension
%---------------------------------------------------------------------------
\usepackage{geometry}
\geometry{
	a4paper,
	left=28mm,
	right=15mm,
	top=30mm,
	headheight=20mm,
	headsep=10mm,
	textheight=242mm,
	footskip=15mm
}
%---------------------------------------------------------------------------

% Makeindex Package
%---------------------------------------------------------------------------
\usepackage{makeidx}                         		% To produce index
\makeindex                                    	% Index-Initialisation
%---------------------------------------------------------------------------

% Intro:
%---------------------------------------------------------------------------
\begin{document}                              	% Start Document
\settocdepth{section}														% Set depth of toc
\pagenumbering{roman}
%---------------------------------------------------------------------------
\providecommand{\titel}{Presentator3000}		%  Hier den Titel des Berichts/Thesis eingeben
\providecommand{\versionnumber}{1.0}			%  Hier die aktuelle Versionsnummer eingeben
\providecommand{\versiondate}{2016-10-11}		%  Hier das Datum der aktuellen Version eingeben

% Set up header and footer
%---------------------------------------------------------------------------
\makeatletter
\patchcmd{\@fancyhead}{\rlap}{\color{bfhgrey}\rlap}{}{}		% new color of header
\patchcmd{\@fancyfoot}{\rlap}{\color{bfhgrey}\rlap}{}{}		% new color of footer
\makeatother

\fancyhf{}																		% clean all fields
\fancypagestyle{plain}{												% new definition of plain style
	\fancyfoot[OR,EL]{\footnotesize \thepage} 	% footer right part --> page number
	\fancyfoot[OL,ER]{\footnotesize \titel, Version \versionnumber, \versiondate}	% footer even page left part
}

\renewcommand{\chaptermark}[1]{\markboth{\thechapter.  #1}{}}
\renewcommand{\headrulewidth}{0pt}				% no header stripline
\renewcommand{\footrulewidth}{0pt} 				% no bottom stripline

\pagestyle{plain}
%---------------------------------------------------------------------------


% Title Page and Abstract
%---------------------------------------------------------------------------
%
% Project documentation template
% ===========================================================================
% This is part of the document "Project documentation template".
% Authors: brd3, kaa1
%

\begin{titlepage}

\graphicspath{{./images/}}


% BFH-Logo absolute placed at (28,12) on A4 and picture (16:9 or 15cm x 8.5cm)
% Actually not a realy satisfactory solution but working.
%---------------------------------------------------------------------------
\setlength{\unitlength}{1mm}
\begin{textblock}{20}[0,0](28,12)
	\includegraphics[scale=1.0]{\detokenize{BFH_Logo_B.png}}
\end{textblock}

\begin{textblock}{154}(28,48)
	\begin{picture}(150,2)
		\put(0,0){\color{bfhgrey}\rule{150mm}{2mm}}
	\end{picture}
\end{textblock}

\begin{textblock}{154}[0,0](28,50)
	\includegraphics[width=15cm]{images/head}			% Titelbild definieren
\end{textblock}

\begin{textblock}{154}(28,134)
	\begin{picture}(150,2)
		\put(0,0){\color{bfhgrey}\rule{150mm}{2mm}}
	\end{picture}
\end{textblock}
\color{black}

% Institution / Titel / Untertitel / Autoren / Experten:
%---------------------------------------------------------------------------
\begin{flushleft}

\vspace*{115mm}

\fontsize{26pt}{28pt}\selectfont
\titel 				\\							% Titel aus der Datei vorspann/titel.tex lesen
\vspace{2mm}

\fontsize{16pt}{20pt}\selectfont\vspace{0.3em}
\vspace{5mm}

\fontsize{10pt}{12pt}\selectfont
\textbf{Report} \\									% eingeben
\vspace{3mm}

% Abstract (eingeben):
%---------------------------------------------------------------------------
\begin{textblock}{150}(28,190)
\fontsize{10pt}{12pt}\selectfont
This document describes the process of what came to be this online presentation tool and what we could have done better in hindsight. \\
\end{textblock}

\begin{textblock}{150}(28,225)
\fontsize{10pt}{17pt}\selectfont
\begin{tabbing}
xxxxxxxxxxxxxxx\=xxxxxxxxxxxxxxxxxxxxxxxxxxxxxxxxxxxxxxxxxxxxxxx \kill
Degree course:	\> Information technology	\\			% Namen eingeben
Lecture:        \> Project 2 \\
Authors:        \> Michael Bolli (michael.bolli.1@students.bfh.ch), \\
\>Robin Schmid (robinsamuel.schmid@students.bfh.ch)		\\					% Namen eingeben
Principal:	    \> Marcel Pfahrer		\\					% Namen eingeben
Date:			\> \versiondate					\\		% aus Datei vorspann/version.tex lesen
\end{tabbing}

\end{textblock}
\end{flushleft}

\begin{textblock}{150}(28,280)
\noindent
\color{bfhgrey}\fontsize{9pt}{10pt}\selectfont
Berner Fachhochschule | Haute école spécialisée bernoise | Bern University of Applied Sciences
\color{black}\selectfont
\end{textblock}


\end{titlepage}

%
% ===========================================================================
% EOF
%
% Versionenkontrolle :
% -----------------------------------------------

\begin{textblock}{180}(15,150)
\color{black}
\begin{huge}
Version history
\end{huge}
\vspace{10mm}

\fontsize{10pt}{18pt}\selectfont
\begin{tabbing}
xxxxxxxxxxx\=xxxxxxxxxxxxxxx\=xxxxxxxxxxxxxx\=xxxxxxxxxxxxxxxxxxxxxxxxxxxxxxxxxxxxxxxxxxxxxxx \kill
Version	\> Date	\> Status		\> Misc.		\\
0.1	\> 2016-10-11	\> Draft		\> Initial version	\\
\end{tabbing}

\end{textblock}

\cleardoubleemptypage
\setcounter{page}{1}
\cleardoublepage
\phantomsection
\let\cleardoublepage\clearpage

%---------------------------------------------------------------------------

% Main part:
%---------------------------------------------------------------------------
\pagenumbering{arabic}


\chapter{Einführung}
\label{chap:intro}

\section{Ausgangslage}
Wenn heute an der Berner Fachhochschule oder anderen Einrichtungen Vorlesungen abgehalten werden, kommen häufig Bildschirmpräsentationen zum Einsatz. Diese unterstützen das Erzählte und helfen den Dozenten, einen roten Faden zu ziehen. Weiter können die Folien (nachfolgend auch Slides genannt) den Unterrichteten abgegeben werden. Diese wiederum können die Unterlagen verwenden, um den Stoff zu repetieren und sich für Examen vorzubereiten. Solche Bildschirmpräsentationen können auf diverse Arten erstellt werden. Die bekannteste Möglichkeit stellt wohl Microsofts PowerPoint dar. Diese Software ist nicht nur sehr anwenderfreundlich sondern auch weit verbreitet. Doch nebst den Anschafungskosten ist PowerPoint abhängig vom Betriebssystem. Für Linux und Unix Systeme (ausser OSX) existiert das Produkt nicht. Gerade in wissenschaftlichen Kreisen verwendet jedoch ein stattlicher Prozentsatz andere Betriebssysteme als Windows/OSX. Als Alternative kommen diverse Varianten zum Einsatz: LibreOffice, OpenOffice, PDFs (aus Latex Vorlagen oder auch andere Quellen), etc. Mit dem aufkommen von Breitband-Internet gibt es auch immer mehr plattformunabhänige Tools, welche direkt im Browser funktionieren. Beispiele dafür sind\emph{Google Slides} oder \emph{Prezi}. 

Nun kommt im spezifischen Fall der BFH TI noch ein weitere Anwendungsfall hinzu: Das Präsentieren von Programmiercode. Dieser Code oder Code-Ausschnitt (nachfolgend auch Code-Fragment, Snippet genannt) ist meist Teil eines grösseren Gesamtkontexts wie etwa einer Übung oder eines Beispiels. Diese Unterlagen werden von den Dozierenden laufend verändert und optimiert. So müssen auch die Präsentationen angepasst werden. Um hier eine Dynamisierung zu schaffen, müsste der Code flexibel in die Präsentation eingebunden werden. Dieser Grundgedanke führte zur Idee eines neuen Produktes, welches nachfolgend erklärt wird.

\section{Idee}
Das Leitmotiv besteht daraus, eine Applikation zu erstellen, welche von Dozenten und Schülern gleichermassen verwendet werden kann. Im Zentrum steht dabei das einfache Erstellen und teilen von Bildschirmpräsentationen. Hinzu kommt das Hinzufügen und Verwalten von Code-Fragmenten. Programmiercode, welcher in einem Repository verfügbar ist, soll auf einfache Weise zu einer Slide hinzugefügt werden können. Wird ein Teil des Snippets geändert, soll die Präsentation mit einer einfachen Aktion (Klick) aktualisiert werden können. Somit gehört das manuelle Anpassen von Code in Präsentationen der Vergangenheit an. 

Um den verschiedenen Geräten und Betriebssystem potenzieller Nutzer Herr zu werden, soll die Applikation als Web-App in gängigen Web-Browsern nutzbar sein.

Eine weitere Funktionalität ist das das Kommentieren von Slides. Hat ein Schüler eine Frage zu einem behandelten Punkt in einem Slide, kann er einen Kommentar hinzufügen. Der Dozierende oder auch Kommilitonen können dann antworten und ihm weiterhelfen. Natürlich kann die Kommentarfunktion auch dazu verwendet werden, um auf Fehler hinzuweisen oder anderen Benutzern weitere Informationen zu verlinken.

Slides sollen nicht fix einer Präsentation angehören, sondern nach Bedarf in eine Präsentation eingefügt werden können. Somit sind Slides wiederverwendbare Komponenten. Um den Überblick an erstellten Slides und Präsentationen zu behalten, lassen sich beide durch das hinzufügen von Tags thematisch kategorisieren. Des weiteren sollen Nutzer sogenannte \emph{Channels} erstellen können, wo sie ihre Präsentationen hinzufügen können. Ein \emph{Channel} dient dazu, anderen Nutzer, beispielsweise Studenten eines Kurses, alle relevanten Präsentationen dieses Kurs geordnet darzustellen. Abonnenten (\emph{Subscriber}) eines \emph{Channels} werden dann benachrichtigt, sobald eine neue Präsentation verfügbar ist oder sich etwas an einer bestehenden Präsentation geändert hat. 

Um die Darstellung von Slides zu beeinflussen sind zwei Funktionen geplant: Mit \emph{Templates} soll eine gewisse Grundstruktur vorgegeben werden. Diese kann in den Editor übernommen und dort bearbeitet werden. Es gibt Templates für Titelfolien, Inhaltsfolien, usw. Die Applikation selbst kommt bereits mit einer Grundmenge an Templates und der User soll nach Bedarf selbst neue hinzufügen können. Um die Gestaltung sind sogenannte \emph{Themes} zuständig. Sie steuern mit CSS-Ausdrücken das Layout einer Präsentation. Auch hier gibt es bereits vorgefertigte Vorlagen, aber auch die Möglichkeit für den User, eigene zu erstellen.

Um den Zuschauern die bestmögliche Erfahrung zu bieten, sollen Präsentationen auch gedruckt beziehungsweise als PDF heruntergeladen werden können. Auch sollen Zuschauer eine Präsentation mit einem spezifischen Link in ihrem eigenen Webbrowser aufrufen können. Dies dient dazu, in überfüllten Hörsälen den Gästen in den hintersten Reihen trotzdem eine optimale Sicht auf die Geschehnisse auf dem Bildschirm zu bieten.

Die aufgezählten Features sind nicht für jeden Benutzer in gleicher Ausführung verfügbar. Je nach Ausgestaltung  des Geschäftsmodel können gewisse Funktionen gar nicht oder nur in limitierter Ausprägung verwendet werden. Es ist aber auch denkbar, das für die Quantität (pro Slide, pro Präsentation, pro Theme, pro Template) eine Microtransaction gestellt wird. Die geeignetste Art der Monetariserung muss in der Entwicklung des Projekts evaluiert werden.



\chapter{Domain Model}
\label{chap:domainmodel}

\begingroup
\setlength{\tabcolsep}{15pt} % Default value: 6pt
\renewcommand{\arraystretch}{1.2}
\begin{table}[h]
\centering
\begin{tabular}{l|c|c|c}
	\textbf{Features} 	&	\textbf{Silver} 	&	\textbf{Gold}	&	\textbf{Diamond} \\ \hline
	Anzahl Präsentationen	&	10	&	unlimitiert	&	unlimitiert \\ \hline
	Anzahl Slides pro Präsentation	&	30	&	unlimitiert	& unlimitiert \\ \hline
	Anzahl Codebases	&	1	& 	unlimitiert	&	unlimitiert \\ \hline
	Channels		&	Nein	&	Nein	& 	Ja, unlimitiert	\\ \hline
	Attachements		&	Nein	&	Nein	&	Ja, unlimitiert \\ \hline
	Preis	&	\textbf{Gratis}	&	\textbf{CHF 5 / Monat} 	& \textbf{CHF 20 / Monat}
\end{tabular}
\end{table}
\endgroup

\begin{figure}[htbp]
    \includesvg[width = 1\textwidth,pdf,path=images/]{erm}
    \caption{Domain Model}
\end{figure}
\chapter{Use Cases}
\label{usecases}

\section{Use Case Diagramm}

\section{Use Case Beschreibungen}
\label{usecase:description}

\subsection{UC 001: Register}
\label{uc:001-register}
\begin{tabular}{|l|p{0.6\textwidth}|}
\hline
\textbf{Use Case Name} 	&	Register	\\ \hline
\textbf{Primary Actor} 	&		\\ \hline
\textbf{Stakeholders}	&		\\ \hline
\textbf{Preconditions}	&		\\ \hline
\textbf{Main Success Scenario}	& 	
\begin{enumerate}
	\item Der EnterpriseManager gibt eine valide Emailadresse und zweimal ein sicheres Passwort ein.
\end{enumerate}
\\ \hline
\textbf{Extensions}	&
\\ \hline
\end{tabular}

\subsection{UC 002: Login}
\label{uc:002-login}

\begin{tabular}{|l|p{0.6\textwidth}|}
\hline
\textbf{Use Case Name} 	&	Login	\\ \hline
\textbf{Primary Actor} 	&	User	\\ \hline
\textbf{Stakeholders}	&	-	\\ \hline
\textbf{Preconditions}	&	Der User ist im System registriert (\nameref{uc:001-register}) und hat seinen Account aktiviert.	\\ \hline
\textbf{Main Success Scenario}	& 	
\begin{enumerate}
	\item Der User gibt seine E-Mail-Adresse und sein Passwort ein
	\item Das System überprüft ob ein Account mit dieser E-Mail Adresse besteht
	\item Das System überprüft, ob das Passwort mit dem des Accounts übereinstimmt
	\item Das System überprüft, ob der Account bereits aktiviert wurde
	\item Der User wird in das System eingeloggt
\end{enumerate}
\\ \hline
\textbf{Extensions}	& 	
\begin{enumerate}
\setcounter{enumi}{1}
    \item
    \begin{enumerate}
    \item Es existiert kein Account mit der eingegebenen E-Mail-Adresse: Das System macht den User darauf aufmerksam
    \end{enumerate}
    \item
    \begin{enumerate}
    \item Das eingegebene Passwort stimmt nicht mit dem Passwort des Accounts überein: Der User wird darauf hingewiesen, dass das Passwort falsch ist und er kann es noch einmal eingeben
    \end{enumerate}
    \item
    \begin{enumerate}
    \item Der Account ist noch nicht aktiviert, was bedeutet, der User hat den Aktivierungslink noch nicht verwendet: Dem User wird angeboten, den Aktivierungslink noch einmal zu versenden
    \end{enumerate}
\end{enumerate}
\\ \hline
\end{tabular}

\subsection{UC 003: Logout}
\label{uc:003-logout}

\begin{tabular}{|l|p{0.6\textwidth}|}
\hline
\textbf{Use Case Name} 	&	Logout	\\ \hline
\textbf{Primary Actor} 	&		\\ \hline
\textbf{Stakeholders}	&		\\ \hline
\textbf{Preconditions}	&	Der User ist im System eingeloggt.(\nameref{uc:002-login}).	\\ \hline
\textbf{Main Success Scenario}	& 	
\begin{enumerate}
	\item Der User wählt die Schaltfläche \emph{Ausloggen} aus
	\item Das System loggt den User aus.
\end{enumerate}
\\ \hline
\textbf{Extensions}	& 	\\ \hline
\end{tabular}

\subsection{UC 004: Manage Presentation}
\label{uc:004-manage-pres}

\begin{tabular}{|l|p{0.6\textwidth}|}
\hline
\textbf{Use Case Name} 	&	Manage Presentation	\\ \hline
\textbf{Primary Actor} 	&		\\ \hline
\textbf{Stakeholders}	&		\\ \hline
\textbf{Preconditions}	&	Anzahl bestehender Präsentationen ist kleiner gleich max_presentations des Abo-Plans.	\\ \hline
\textbf{Main Success Scenario}	&
\begin{enumerate}
	\item Ein eingeloggter User kann Präsentationen hinzufügen
	\item Eine Präsentation besteht aus Slides in einer definierten Reihenfolge
	\item Eine Präsentation ist in der Grösse durch max_slides_pp des Abo-Plans beschränkt.
\end{enumerate}
\\ \hline
\textbf{Extensions}	& 	\\ \hline
\end{tabular}

\subsection{UC 005: Manage Slide}
\label{uc:005-manage-slide}

\begin{tabular}{|l|p{0.6\textwidth}|}
\hline
\textbf{Use Case Name} 	&	Manage Slide	\\ \hline
\textbf{Primary Actor} 	&		\\ \hline
\textbf{Stakeholders}	&		\\ \hline
\textbf{Preconditions}	&	Anzahl Slides in einer Präsentation ist kleiner gleich max_slides_pp des Abo-Plans.	\\ \hline
\textbf{Main Success Scenario}	&
\begin{enumerate}
	\item Ein eingeloggter User kann einer Präsentation Slides hinzufügen
	\item Ein Slide kann in einer anderen Präsentation wiederverwendet werden (shared = true)
	\item Ein Slide kann mithilfe eines WYSIWYG-Editors benutzerfreundlich bearbeitet werden
	\item Slides können in der Reihenfolge via Drag'n'Drop angepasst werden
\end{enumerate}
\\ \hline
\textbf{Extensions}	& 	\\ \hline
\end{tabular}

\subsection{UC 006: Share Presentation}
\label{uc:006-share-pres}

\subsection{UC 007: View Presentation}
\label{uc:007-view-pres}

\begin{tabular}{|l|p{0.6\textwidth}|}
\hline
\textbf{Use Case Name} 	&	View Presentation	\\ \hline
\textbf{Primary Actor} 	&		\\ \hline
\textbf{Stakeholders}	&		\\ \hline
\textbf{Preconditions}	&		\\ \hline
\textbf{Main Success Scenario}	&
\begin{enumerate}
	\item Die Präsentation wird angezeigt
\end{enumerate}
\\ \hline
\textbf{Extensions}	& 	\\ \hline
\end{tabular}

\subsection{UC 008: Print Presentation}
\label{uc:008-print-pres}

\subsection{UC 009: Manage Comment}
\label{uc:009-manage-comment}

\subsection{UC 010: Manage Codebase}
\label{uc:010-manage-codebase}

\begin{tabular}{|l|p{0.6\textwidth}|}
\hline
\textbf{Use Case Name} 	&	Manage Codebase	\\ \hline
\textbf{Primary Actor} 	&		\\ \hline
\textbf{Stakeholders}	&		\\ \hline
\textbf{Preconditions}	&	Anzahl bestehender Codebases ist kleiner gleich max_codebases des Abo-Plans.	\\ \hline
\textbf{Main Success Scenario}	&
\begin{enumerate}
	\item Ein eingeloggter User kann Codebases hinzufügen, bspw. Github
	\item Eine Codebase hat eine URL zu einem Git-Repo und einen Titel
\end{enumerate}
\\ \hline
\textbf{Extensions}	& 	\\ \hline
\end{tabular}

\subsection{UC 011: Manage Code Fragment}
\label{uc:011-manage-code-frag}

\begin{tabular}{|l|p{0.6\textwidth}|}
\hline
\textbf{Use Case Name} 	&	Manage Code Fragment	\\ \hline
\textbf{Primary Actor} 	&		\\ \hline
\textbf{Stakeholders}	&		\\ \hline
\textbf{Preconditions}	&	Eine Codebase existiert	\\ \hline
\textbf{Main Success Scenario}	&
\begin{enumerate}
	\item Bei Erstellen eines Slides kann Code direkt aus einer Codebase angezeigt werden
	\item Dazu muss im WYSIWYG-Editor der entsprechende Button gewählt werden. Folgende Felder müssen ausgefüllt werden:
	\begin{enumerate}
	    \item Filepath
	    \item Revision\#
	    \item Branch
	    \item Start [Line]
	    \item Ende [Line]
	\end{enumerate}
	\item Das System generiert so im Hintergrund einen Shortcode, der auch im Source-Modus des Editors manuell eingegeben werden kann
\end{enumerate}
\\ \hline
\textbf{Extensions}	& 	\\ \hline
\end{tabular}

\subsection{UC 012: Update External Resources}
\label{uc:012-update-external-res}

\begin{tabular}{|l|p{0.6\textwidth}|}
\hline
\textbf{Use Case Name} 	&	Update External Resources	\\ \hline
\textbf{Primary Actor} 	&		\\ \hline
\textbf{Stakeholders}	&		\\ \hline
\textbf{Preconditions}	&		\\ \hline
\textbf{Main Success Scenario}	&
\begin{enumerate}
	\item Die Präsentation wird angezeigt
\end{enumerate}
\\ \hline
\textbf{Extensions}	& 	\\ \hline
\end{tabular}

\subsection{UC 013: Manage Channel}
\label{uc:013-manage-channel}

\begin{tabular}{|l|p{0.6\textwidth}|}
\hline
\textbf{Use Case Name} 	&	Manage Channel	\\ \hline
\textbf{Primary Actor} 	&	User	\\ \hline
\textbf{Stakeholders}	&	-	\\ \hline
\textbf{Preconditions}	&	Der Abo-Plan des Users hat enable_channels auf true.	\\ \hline
\textbf{Main Success Scenario}	&
\begin{enumerate}
	\item Ein angemeldeter User kann Channels erstellen
	\item In diesen können Präsentationen gruppiert und geordnet werden.
\end{enumerate}
\\ \hline
\textbf{Extensions}	& 	\\ \hline
\end{tabular}

\subsection{UC 014: Subscribe Channel}
\label{uc:014-subscribe-channel}

\begin{tabular}{|l|p{0.6\textwidth}|}
\hline
\textbf{Use Case Name} 	&	Subscribe Channel	\\ \hline
\textbf{Primary Actor} 	&		\\ \hline
\textbf{Stakeholders}	&		\\ \hline
\textbf{Preconditions}	&		\\ \hline
\textbf{Main Success Scenario}	&
\begin{enumerate}
	\item Ein angemeldeter User kann Channels abonnieren.
	\item Abonnierte Channels tauchen auf seiner Homepage auf.
	\item Bei neuen Präsentationen in einem abonnierten Channel wird eine E-Mail Benachrichtigung ausgelöst.
\end{enumerate}
\\ \hline
\textbf{Extensions}	& 	\\ \hline
\end{tabular}

\subsection{UC 015: Bookmark Presentation}
\label{uc:015-bookmark-pres}

\subsection{UC 016: Manage Profile}
\label{uc:016-manage-profile}

\chapter{Solution}
\label{chap:solution}
After some fiddling in the beginning, implementation was quite straightforward. The grammar had to be converted into a tree structure, each node representing methods in the Logo.jj file. Most of these methods are writing strings of Java code into the destination file via PrintWriter. It was of course necessary to pay attention to things like indentation and REPCOUNT in- and decreasing.

The REPCOUNT (or for-loop in Java) isn't directly supported by LOGO. We solved that with a similar approach to indentation: An integer keeps track of REPCOUNT, which is in- and decreased when used in tokens like REPEAT.

\begingroup
\renewcommand{\cleardoublepage}{}
\renewcommand{\clearpage}{}
\chapter{Test}
\label{chap:test}
Apart from the provided LOGO files in the logofiles Folder, we used our own test file test.logo for final testing.
\endgroup

\begingroup
\renewcommand{\cleardoublepage}{}
\renewcommand{\clearpage}{}

\chapter{Limitations \& Considerations}
\label{chap:limits}
This parser only verifies the correctness of the LOGO program on a syntactical level. The semantic verficiation will be performed during the Java compilation of what is produced by our parser.

We might have been faster to solve this project if our group had met in real life for a day, instead of communicating via e-Mail over weeks.
\endgroup
\end{document}
